\documentclass[../main.tex]{subfiles}
\begin{document}

\chapter*{Kết luận}
\addcontentsline{toc}{chapter}{Kết luận}

Trong quá trình thực hiện đồ án, em đã có cơ hội tìm hiểu sâu sắc về một trong những mô hình quan trọng nhất của lĩnh vực Tối ưu hóa: Bài toán vận tải. Qua việc nghiên cứu cơ sở lý thuyết \cite{hitchcock1941, koopmans1949}, từ việc thiết lập mô hình toán học ban đầu, tìm phương án cực biên xuất phát bằng các phương pháp như góc Tây Bắc, cực tiểu chi phí và xấp xỉ Vogel, cho đến việc ứng dụng thuật toán thế vị của Dantzig \cite{dantzig1951} để tối ưu hóa chi phí vận chuyển, em đã nắm vững được bản chất và quy trình giải quyết trọn vẹn của bài toán này theo hệ thống kiến thức giáo trình chuẩn \cite{bachkim}.

Bên cạnh việc củng cố lý thuyết toán học, đồ án còn giúp em rèn luyện tư duy logic và kỹ năng lập trình thực tế thông qua việc tự tay xây dựng một phần mềm giải toán vận tải dưới nền tảng ứng dụng web. Ứng dụng không chỉ có khả năng tính toán chuẩn xác và tự động bao quát mọi kịch bản hóc búa của bài toán (như bài toán không cân bằng thu phát, hiện tượng suy biến sớm hay xuất hiện tập phương án tối ưu tổng quát), mà còn được thiết kế với một giao diện cực kỳ trực quan, thân thiện. Việc có thể mô phỏng tường minh từng bước lặp thay đổi các giá trị $u_i, v_j, \Delta_{ij}$ và biểu diễn chu trình chuyển hàng trực tiếp trên bảng vận tải giúp cho mọi người dùng đều có thể dễ dàng theo dõi, hiểu rõ cơ chế của thuật toán thế vị nhanh chóng và sâu sắc hơn rất nhiều so với cách giải tay truyền thống.

Mặc dù đồ án đã hoàn thành tốt những mục tiêu cơ bản đề ra ban đầu, phần mềm vẫn còn một số khía cạnh có thể tiếp tục được phát triển mở rộng trong tương lai, chẳng hạn như: tối ưu hiệu suất xử lý bằng cấu trúc dữ liệu tốt hơn để giải các bài toán với kích thước hàng chục ngàn trạm quy mô lớn, hay bổ sung thêm tính năng nhập/xuất dữ liệu ma trận qua tệp tin Excel nhằm tăng cường tính ứng dụng thực tiễn trong công nghiệp. Em hy vọng đồ án này sẽ trở thành một bước đệm vững chắc giúp bản thân có thêm nhiều động lực để tiếp tục nghiên cứu, ứng dụng các công cụ tin học và thuật toán tối ưu vào việc tự động hóa giải quyết các bài toán lớn hơn trong thực tế.

Em xin gửi lời cảm ơn chân thành đến các Thầy Cô giáo trong Bộ môn đã luôn tâm huyết truyền đạt kiến thức, hướng dẫn và truyền cảm hứng cho em trong suốt quá trình học tập và hoàn thiện đồ án này.

\end{document}
